\specsection{Inverse limits of measures}

This is a theory that has developed mainly in response to the needs of Probability Theory. Problems concerning a finite sequence of random variables \(X_{1},\ldots,X_{n}\) are in principle solved when one knows the law \(P_{X}\) of the sequence: this is a positive measure on \(R^{n}\) of mass 1 such that the probability of obtaining simultaneously the inequalities \(a_{1}\leqslant X_{1}\leqslant b_{1},\ldots,a_{n}\leqslant X_{n}\leqslant b_{n}\) is equal to \(\mathrm{P}_{\mathrm{X}}(\mathrm{C})\) , where C is the closed box \([a_{1},b_{1}]\times\cdots\times[a_{n},b_{n}]\) in \(R^{n}\) . In practice, the measure \(P_{X}\) either has discrete support or admits a density with respect to Lebesgue measure. When dealing with an infinite sequence \((\mathrm{X}_{n})_{n\geqslant1}\) of random variables, one generally knows the law \(P_{n}\) of the partial sequence \((\mathrm{X}_{1},\ldots,\mathrm{X}_{n})\) for every integer \(n\geqslant1\) ; these data satisfy a compatibility condition that expresses that the sequence \((\mathrm{P}_{n})_{n\geqslant1}\) is an inverse system (or `projective system') of measures. Until about 1920, the probabilities of events relating to the infinite sequence were defined in more or less implicit fashion by `natural' passages to the limit from the probabilities of the finite case; one thus assumes that the probability that a game terminates is the limit, as n tends to infinity, of the probability that it terminates in at most n steps. Naturally, such a theory is not very coherent, and nothing excludes the presence of `paradoxes', the same probability receiving two distinct estimates according as one evaluates by means of one or the other of two procedures, each as `natural' as the other.

Steinhaus (V) appears to have been the first to have felt the need for considering (for the game of heads-or-tails) not only the inverse system \((\mathbf{P}_{n})_{n\geqslant1}\) but its limit as well. A little earlier, in 1919, Daniell (VI, b)) had proved in general the existence of such inverse limits, \(^{(7)}\) but this result seems to have remained unknown in Europe. It was rediscovered in 1933 by Kolmogoroff in the work (XII) where the author formulates the axiomatic conception of Probability Theory. The proofs of Daniell and Kolmogoroff make use of a compactness argument, which is substantially the same as the one we have employed in Th. 2 of §4, No.~3 and which rests on Dini's theorem.

The Daniell--Kolmogoroff theorem left nothing to be desired for the case of random sequences \((\mathbf{X}_{n})_{n\geqslant1}\) , but the study of random functions undertaken from 1935 on by Kolmogoroff, Feller and Doob harbors difficulties of quite another order. Consider for example an interval T of R, representing the set of instants of observation of a `stochastic process'; the set of possible trajectories is the product space \(R^{T}\) , regarded as the inverse limit of the partial products \(R^{H}\) , where H runs over the set of finite subsets of T; one generally assumes given an inverse system of measures \((\mu_{\mathrm{H}})\) (cf.~§4, No.~2). Kolmogoroff's theorem does indeed yield a measure on \(R^{T}\) , but it is defined on a tribe notably smaller than the Borel tribe. \(^{(8)}\) A variant of Kolmogoroff's construction, which yields a measure on a topological space, is due to Kakutani (Proc. Imp. Acad. Tokyo 19 (1943), 184--188), and has been rediscovered several times since: one regards \(\mu_{H}\) as a measure on \(\overline{R}^{H}\) carried by \(R^{H}\) ; \(^{(9)}\) the compact space \(E = \overline{R}^{T}\) is the inverse limit of the finite products \(\overline{R}^{H}\) and one can define a measure \(\mu\) on E as the inverse limit of the \(\mu_{H}\) (cf.~Ch. III, §4, No.~5). However, this procedure has a serious inconvenience; the elements of \(\overline{R}^{T}\) possess no regularity property that permits advancing the probabilistic study of the process---or even simply eliminating the parasitic values \(\pm\infty\) introduced by the compactification \(\overline{R}\) of R. This can be remedied by inducing the measure \(\mu\) of \(\overline{R}^{T}\) on a particular subspace (for example \(\mathcal{C}(T)\) in the case of Brownian motion); the fundamental difficulty arises from the fact that a function space, even one of the usual type, is not necessarily \(\mu\) -measurable in \(\overline{R}^{T}\) , and even the choice of function space may be questionable. \(^{(10)}\)

A decisive step was taken in 1956 by Prokhorov in a work (XIII) that has had a determining influence on the theory of stochastic processes. By putting into a suitable axiomatic form the methods used by Wiener in the article analyzed above, he established a general existence theorem for inverse limits of measures on function spaces that is a special case of Th. 1 of §4, No.~2 corresponding to Polish spaces.

A more restricted class of inverse systems was introduced by Bochner (XIV) in 1947; these are inverse systems formed by finite-dimensional real vector spaces and surjective linear mappings. The inverse limit of such a system may be identified in a natural way with the algebraic dual \(\mathbf{E}^*\) of a real vector space \(\mathbf{E}\) , equipped with the weak topology \(\sigma (\mathbf{E}^{*},\mathbf{E})\) ; a corresponding inverse system of measures has a limit that is a measure \(\mu\) defined on a tribe notably smaller than the Borel tribe of \(\mathbf{E}^*\) . Bochner completely characterized such `promasures' by their Fourier transform, which is a function on \(\mathbf{E}\) . But this result is scarcely usable in the absence of a topology on \(\mathbf{E}\) , in which case one has to examine the possibility of regarding \(\mu\) as a measure on the topological dual \(\mathbf{E}'\) of \(\mathbf{E}\) . In an independent way, R. Fortet and E. Mourier, while seeking to generalize to random variables with values in a Banach space certain classical results of Probability Theory (law of large numbers, central limit theorem) also brought into evidence the fundamental role played by the Fourier transformation in such questions. But substantial progress was not made until 1956 when Gelfand (XV, \(b\) ) suggested that the natural setting for the Fourier transformation is not that of Banach spaces or Hilbert spaces, but that of nuclear Fréchet spaces. He conjectured that every continuous function of positive type on such a space is the Fourier transform of a measure on its dual, a result established soon afterwards by Minlos (XVI). Its importance stems above all from the fact that it is applicable to spaces of distributions, and that the quasi-totality of function spaces are Borel subsets of the space of distributions (which thus constitutes a much better receptacle than \(\mathbf{R}^{\mathrm{T}}\) ). \(^{(11)}\) The theory of random distributions is an area in full expansion, and we shall content ourselves by referring the reader to the book of Gelfand and Vilenkin (XVII).

The results on inverse limits just mentioned make use of the existence of topologies on the base spaces. One may ask whether there exists an analogous theory in the case of `abstract' measures. Von Neumann proved in 1935 the existence of product measures in all cases, but the discovery of a counter-example by Jessen and Andersen (XVIII) dashed the hope that every inverse system of measures admits a limit. Two palliatives have been discovered: in 1949, C. Ionescu-Tulcea established the existence of countable inverse limits, by means of the existence of suitable disintegrations, \(^{(12)}\) a highly interesting result for the study of Markoff processes; moreover, it had been observed that the topology of the spaces only intervenes through the intermediary of the set of compact subsets. It was therefore natural to try to axiomatize this situation within the abstract theory, by means of the concept of compact class of subsets of a set. This work was done in 1953 by Marczewski (who established an abstract inverse limit theorem by such means) and Ryll-Nardzewski (who treated the disintegration of measures). \(^{(13)}\)

